% !TeX program = lualatex
% !TeX encoding = utf8
% !TeX spellcheck = uk_UA
% !TeX root = ../main.tex

%% --------------------------------------------------------
%% Chapter file: LinearWaves.tex
%% --------------------------------------------------------

\graphicspath{{\currfilebase/Pictures}}



\chapter{Хвилі в~плазмі}
\label{sec:plasma-waves}\hypertarget{plasma-waves}{}
%% --------------------------------------------------------

Хвиля~--- це не обов'язково перенесення речовини. У~найпростішому розумінні хвиля є
поширенням збурення, яке завдяки внутрішнім силам системи передається від однієї її
ділянки до іншої.

Для плазми це твердження має особливо багатий зміст. Плазма одночасно має тиск, електричні
заряди, струми та магнітне поле. Тому після малого збурення відновлювальну дію можуть
створювати зовсім різні фізичні механізми. У~найпростішому випадку плазму повертає до
рівноваги градієнт тиску; у~замагніченій плазмі додатково виникає натяг силових ліній
магнітного поля.

Саме тому в~плазмі існує не одна \enquote{швидкість звуку}, а~ціла система хвильових мод.



%% ========================================================
\section{Що в~плазмі може бути відновлювальною силою?}
%% ========================================================



Корисно почати не з~дисперсійного рівняння, а~з~фізичної картини. Уявімо, що деяку ділянку
плазми трохи змістили від положення рівноваги. Щоб виникла хвиля, це зміщення повинно
породити силу, яка намагається повернути систему назад.

У~звичайному газі такою силою є градієнт тиску. Якщо невеликий об'єм газу стиснути, його
тиск зростає і~виникає сила, спрямована назовні. Якщо ж цей об'єм розрідити, тиск падає,
і~сусідні області штовхають його назад.

У~плазмі до цього додається магнітне поле. Викривлення силових ліній породжує магнітний
натяг, а~зміна величини поля~--- зміну магнітного тиску

\begin{equation}
p_B=\frac{B^2}{8\pi}.
\label{eq:wave-magnetic-pressure}
\end{equation}

Тому вже на якісному рівні можна передбачити три важливі типи відновлювальної дії:

\begin{enumerate}
\item газодинамічний тиск;
\item натяг магнітних силових ліній;
\item зміна магнітного тиску.
\end{enumerate}

Перший механізм дає звукову хвилю, другий~--- альвенівську, а поєднання тиску плазми та
магнітного поля породжує магнітозвукові хвилі.

\begin{figure}[ht]
\centering
\fbox{%
  \parbox[c][62mm][c]{0.82\textwidth}{%
    \centering
    \textbf{ЗАГЛУШКА ДЛЯ РИСУНКА}\\[2mm]
    Одна й~та сама рівноважна плазма та три способи її збурити.\\[2mm]
    \textbf{Ліворуч:} стиснення та розрідження~--- градієнт тиску.\\
    \textbf{По центру:} поперечне зміщення силових ліній --- магнітний натяг.\\
    \textbf{Праворуч:} стискання плазми разом зі зміною $B$ --- магнітний тиск + газовий тиск. }%
}
\caption{Фізичне походження основних МГД-хвиль: різні відновлювальні сили породжують різні
способи коливання плазми.}
\label{fig:plasma-wave-restoring-forces}
\end{figure}

Ця картина буде нашим орієнтиром у~всьому розділі. Спочатку розглянемо кожен механізм
окремо, а~потім побачимо, як вони змішуються в магнітованій плазмі.



%% ========================================================
\section{Звукова хвиля}
%% ========================================================



Почнемо з~найпростішого випадку~--- плазми без магнітного поля. Якщо плазма перебуває
в~локальній термодинамічній рівновазі, її можна розглядати як звичайне стисливе
середовище.

Рівняння неперервності має вигляд

\begin{equation}
\frac{\partial \rho_m}{\partial t}
+\nabla\cdot(\rho_m\vect v)=0.
\label{eq:wave-continuity}
\end{equation}

Для малих швидкостей та збурень однорідної рівноваги $\rho_m=\rho_0+\delta\rho$ рівняння
\eqref{eq:wave-continuity} перетворюється на

\begin{equation}
\frac{\partial\delta\rho}{\partial t}
+\rho_0\nabla\cdot\delta\vect v=0.
\label{eq:wave-linear-continuity}
\end{equation}

Рівняння руху в~цьому наближенні має вигляд

\begin{equation}
\rho_0\frac{\partial\delta\vect v}{\partial t}
=-\nabla\delta p.
\label{eq:wave-linear-momentum}
\end{equation}

Щоб замкнути систему, потрібен зв'язок між збуреннями тиску та густини. Для адіабатичної
хвилі

\begin{equation}
\delta p=c_s^2\delta\rho,
\qquad
c_s^2=\left(\frac{\partial p}{\partial\rho_m}\right)_S.
\label{eq:wave-sound-speed}
\end{equation}

Тут $c_s$~--- швидкість звуку.

Підставляючи \eqref{eq:wave-sound-speed} у \eqref{eq:wave-linear-momentum}, а~потім
використовуючи \eqref{eq:wave-linear-continuity}, одержуємо хвильове рівняння

\begin{equation}
\frac{\partial^2\delta\rho}{\partial t^2}
-c_s^2\nabla^2\delta\rho=0.
\label{eq:wave-sound-equation}
\end{equation}

Отже, звук виникає тому, що локальне стискання породжує зміну тиску, а~зміна тиску
породжує рух, який передає збурення сусіднім ділянкам.

Для плазми це лише перша сходинка. На відміну від нейтрального газу, плазмовий тиск може
бути пов'язаний із температурою електронів та іонів, а~на малих масштабах до руху частинок
підключаються електричні поля. Тому кінетичний опис плазми дає вже інші, зокрема
іонно-звукові, хвилі.

\begin{figure}[ht]
\centering
\fbox{%
  \parbox[c][55mm][c]{0.82\textwidth}{%
    \centering
    \textbf{ЗАГЛУШКА ДЛЯ РИСУНКА}\\[2mm]
    Поздовжня звукова хвиля в~плазмі.\\
    Показати області стискання та розрідження, напрямок швидкості частинок і~градієнт
    тиску.\\[2mm]
    Бажано показати, що частинки коливаються переважно вздовж напрямку поширення хвилі, а~не \enquote{їдуть} разом із хвилею. }%
}
\caption{Звукова хвиля як поширення чергувань стискання та розрідження.}
\label{fig:plasma-sound-wave}
\end{figure}



%% ========================================================
\section{Альвенівська хвиля}
%% ========================================================



Тепер увімкнемо однорідне магнітне поле $\vect B_0$.

З'являється принципово новий ступінь свободи. Силову лінію магнітного поля можна
викривити. Але викривлена силова лінія має натяг і~прагне повернутися до прямолінійного
стану.

Це можна уявити як струну: якщо її відхилити поперек початкового напрямку, натяг струни
породжує силу, яка намагається її випрямити. Для магнітного поля роль натягу визначається
другим доданком у розкладі магнітної сили:

\begin{equation}
\frac{1}{c}[\vect j\times\Bfield]
=
-\nabla\frac{B^2}{8\pi}
+
\frac{1}{4\pi}(\Bfield\cdot\nabla)\Bfield.
\label{eq:wave-magnetic-force}
\end{equation}

Для поперечного збурення, що поширюється вздовж $\vect B_0$, основною відновлювальною
силою є саме магнітний натяг.

Лінеаризовані рівняння ідеальної МГД приводять до хвильового рівняння

\begin{equation}
\frac{\partial^2\delta\vect v_\perp}{\partial t^2}
=
v_A^2
\frac{\partial^2\delta\vect v_\perp}{\partial z^2},
\label{eq:wave-alfven-equation}
\end{equation}

де

\begin{equation}
v_A=\frac{B_0}{\sqrt{4\pi\rho_0}}
\label{eq:wave-alfven-speed}
\end{equation}

--- швидкість Альвена.

Для гармонічної хвилі $\delta\vect v_\perp\propto \exp[i(kz-\omega t)]$ маємо

\begin{equation}
\omega^2=k^2v_A^2,
\qquad
\omega=\pm kv_A.
\label{eq:wave-alfven-dispersion}
\end{equation}

Отже, альвенівська хвиля є поперечною: плазма та магнітне поле збурюються поперек напрямку
поширення, тоді як саме поле забезпечує відновлювальну силу.

\begin{figure}[ht]
\centering
\fbox{%
  \parbox[c][68mm][c]{0.82\textwidth}{%
    \centering
    \textbf{ЗАГЛУШКА ДЛЯ РИСУНКА}\\[2mm]
    Альвенівська хвиля вздовж однорідного поля $\vect B_0$.\\[2mm]
    Показати майже прямі силові лінії, які послідовно зміщуються поперек напрямку $z$. На
    них накласти $\delta\vect B_\perp$ та $\delta\vect v_\perp$.\\[2mm]
    Окремою стрілкою показати напрямок поширення хвилі. }%
}
\caption{Альвенівська хвиля: поперечне збурення поширюється вздовж магнітних силових ліній
завдяки їхньому натягу.}
\label{fig:alfven-wave}
\end{figure}

Тут уже видно важливу відмінність від звуку. У~звуковій хвилі відновлювальним механізмом
був тиск. В~альвенівській хвилі ним є магнітний натяг. Тому швидкість Альвена визначається
не температурою, а~співвідношенням магнітного поля та масової густини, як видно з
\eqref{eq:wave-alfven-speed}.



%% ========================================================
\section{Магнітозвукова хвиля}
%% ========================================================



У~реальній магнітованій плазмі магнітний натяг не є єдиною відновлювальною силою. Якщо
збурення стискає плазму, воно змінює не лише густину та тиск, а~й~величину магнітного
поля.

Тому виникає зв'язок між двома механізмами:

\begin{equation*}
\text{газовий тиск}
\qquad\Longleftrightarrow\qquad
\text{магнітний тиск}.
\end{equation*}

Результатом є магнітозвукові хвилі.

Для простоти розглянемо поширення вздовж напрямку, перпендикулярного до однорідного
магнітного поля. Лінеаризація МГД-рівнянь приводить до дисперсійного співвідношення

\begin{equation}
\omega^2=k^2(c_s^2+v_A^2).
\label{eq:wave-fast-magnetosonic}
\end{equation}

Відповідна швидкість

\begin{equation}
v_{\mathrm{ms}}=\sqrt{c_s^2+v_A^2}.
\label{eq:wave-magnetosonic-speed}
\end{equation}

Ця формула добре демонструє фізичну ідею. Дві незалежні відновлювальні дії працюють
одночасно: тиск плазми дає внесок $c_s^2$, а~магнітне поле~--- внесок $v_A^2$.

\begin{figure}[ht]
\centering
\fbox{%
  \parbox[c][65mm][c]{0.82\textwidth}{%
    \centering
    \textbf{ЗАГЛУШКА ДЛЯ РИСУНКА}\\[2mm]
    Магнітозвукова хвиля, що поширюється поперек $\vect B_0$.\\[2mm]
    Показати області стискання плазми та силових ліній: де плазма стискається, там зростає
    $B$ і~магнітний тиск.\\[2mm]
    Підписати два внески у~відновлювальну силу: \enquote{тиск плазми} та \enquote{магнітний тиск}. }%
}
\caption{Магнітозвукова хвиля: стискання плазми супроводжується зміною магнітного поля.}
\label{fig:magnetosonic-wave}
\end{figure}

У~загальному випадку напрямок поширення не зобов'язаний бути ні паралельним, ні
перпендикулярним до магнітного поля. Тоді отримуємо дві магнітозвукові моди: швидку та
повільну.

Їхні фазові швидкості визначаються співвідношенням

\begin{equation}
v_{\pm}^2
=
\frac{1}{2}
\left[
v_A^2+c_s^2
\pm
\sqrt{
(v_A^2+c_s^2)^2
-4v_A^2c_s^2\cos^2\theta
}
\right],
\label{eq:wave-magnetosonic-general}
\end{equation}

де $\theta$~--- кут між $\vect k$ та $\vect B_0$.

Знак \enquote{$+$} відповідає швидкій моді, а~знак \enquote{$-$}~--- повільній.

Формула \eqref{eq:wave-magnetosonic-general} на перший погляд виглядає складною. Її
фізичний зміст, однак, простий: при довільному напрямку поширення рух плазми може
одночасно викривляти поле, стискати його та змінювати густину. Тому \enquote{чиста}
звукова або \enquote{чиста} альвенівська картина більше не є достатньою.



%% ========================================================
\section{Три основні МГД-хвилі}
%% ========================================================



Тепер можна звести картину в~одну схему.

\begin{table}[ht]
\centering
\begin{tblr}{
  colspec={X[l,m]X[l,m]X[l,m]X[c,m]},
  row{odd}={gray!10},
  hlines,
  vlines
}
Хвиля & Основна відновлювальна сила & Поляризація & Швидкість \\

Звукова &
градієнт тиску &
поздовжня &
$c_s$ \\

Альвенівська &
магнітний натяг &
поперечна &
$v_A$ \\

Швидка магнітозвукова &
тиск + магнітне поле &
стискувальна &
$v_+$ \\

Повільна магнітозвукова &
тиск + магнітне поле &
змішана &
$v_-$
\end{tblr}
\caption{Основні хвильові моди в~рідинному описі плазми.}
\label{tab:plasma-wave-modes}
\end{table}

Таблиця~\ref{tab:plasma-wave-modes} дає корисну фізичну карту, але її не слід сприймати як
остаточну класифікацію всіх плазмових хвиль. Ми розглядали тут лише МГД-моди~--- тобто
довгохвильові та повільні процеси, для яких рідинний опис є достатнім.

На менших масштабах стають важливими окремі рухи частинок. Тоді виникають кінетичні моди,
для яких поняття \enquote{швидкість звуку} або \enquote{натяг силових ліній} вже не дає
повної картини.



%% ========================================================
\section{Що змінюється на менших масштабах?}
%% ========================================================



До цього моменту ми поводилися з~плазмою як із суцільним середовищем. Це означає, що
інформація про окремі частинки була усереднена і замість функції розподілу
використовувалися густина, швидкість, тиск та інші моменти.

Таке усереднення має межі.

Якщо довжина хвилі стає порівнянною з~ларморівським радіусом або довжиною вільного
пробігу, структура руху частинок уже не може бути повністю прихована всередині
коефіцієнтів МГД. Аналогічно, якщо частота хвилі наближається до характерних частот руху
частинок, виникають резонансні ефекти.

Тоді треба повернутися до функції розподілу

\begin{equation}
f_s=f_s(\vect r,\vect v,t)
\end{equation}

і~досліджувати не лише рух середнього потоку, а~й~те, як хвиля взаємодіє з~частинками
різних швидкостей.

Саме тут починається кінетична теорія хвиль.

У~наступному розділі ми тому не просто введемо ще один спосіб записувати рівняння. Ми
змінимо рівень опису: замість суцільного середовища розглядатимемо розподіл частинок
у~фазовому просторі. Це дозволить пояснити, звідки беруться кінетичні хвилі, дисперсія
Ленгмюра та загасання Ландау.


% --------------------------------------------------------
\subsection{Від МГД-хвиль до кінетичного опису}
% --------------------------------------------------------

У~попередніх розділах ми розглядали плазму як суцільне середовище. Такий опис дозволив
нам отримати три основні типи хвиль в~ідеальній однорідній МГД: альвенівську,
швидку магнітозвукову та повільну магнітозвукову хвилі. Усі вони були отримані з~рівнянь
для макроскопічних величин --- густини, швидкості потоку, тиску та магнітного поля.

Виникає природне питання: \emph{наскільки фундаментальним є такий опис?}

Адже плазма насправді складається з~окремих заряджених частинок. Ми вже знаємо, що
їхній стан описує функція розподілу $f_s(\vect r,\vect v,t)$, тоді як МГД працює лише
з~кількома її моментами. Отже, між кінетичною та МГД-картинами має існувати певний
зв'язок.

Цей зв'язок можна уявити так:

\begin{equation*}
\tcbhighmath{
\text{функція розподілу}
;\longrightarrow;
\text{моменти}
;\longrightarrow;
\text{рідинні рівняння}
;\longrightarrow;
\text{МГД}
}
\end{equation*}

Перехід відбувається не тому, що частинки перестають бути важливими. Навпаки, їхня
мікроскопічна динаміка залишається тією самою; ми лише відмовляємося від детальної
інформації про розподіл частинок за швидкостями, якщо вона не потрібна для задачі.

Те саме стосується хвиль. Кінетична теорія допускає набагато ширший клас колективних
рухів, а~МГД-хвилі відповідають певній довгохвильовій та низькочастотній межі цього
опису.

\begin{figure}[ht]
\centering
\begin{forest}
for tree={
draw,
rounded corners,
align=center,
edge={->},
l sep=12mm,
s sep=7mm,
font=\small,
}
[Кінетичний опис
[Рідинний опис
[Магнітогідродинаміка
[Альвенівська хвиля]
[Швидка магнітозвукова хвиля]
[Повільна магнітозвукова хвиля]
]
]
[Кінетичні поправки
[скінченний ларморівський радіус]
[резонансна взаємодія]
[залежність від розподілу за швидкостями]
]
]
\end{forest}
\caption{Схематичний зв'язок між кінетичним і~МГД-описами хвиль. МГД-хвилі є
великомасштабною межею загальної кінетичної картини.}
\label{fig:kinetic-mhd-wave-bridge}
\end{figure}

% --------------------------------------------------------
\subsubsection{Що означає ``кінетичне продовження'' МГД-хвилі?}
% --------------------------------------------------------

Розглянемо, наприклад, альвенівську хвилю. У~межах ідеальної МГД для однорідної плазми
її дисперсійне співвідношення має вигляд

\begin{equation}\label{eq:mhd-alfven-bridge}
\omega=k_\parallel v_A.
\end{equation}

У~цьому результаті немає явної залежності від швидкості окремої частинки. Це природно:
МГД описує не індивідуальні траєкторії, а~середній рух плазми як суцільного середовища.

Тепер уявімо, що ми описуємо ту саму плазму не через густину, швидкість і~магнітне поле,
а~через функції розподілу $f_s$. У~такому описі альвенівська хвиля не зникає. Вона
повинна з'явитися як один із~розв'язків кінетичних рівнянь у~відповідній межі параметрів.

Отже, кінетична теорія не заперечує результат \eqref{eq:mhd-alfven-bridge}. Навпаки,
у~режимі, де МГД застосовна, кінетичний опис має відтворити цей результат.

Саме тому правильніше говорити не про дві різні хвилі, а~про \emph{два рівні опису
одного й~того самого колективного руху}.

У~МГД ми запитуємо:

\begin{equation*}
\text{як змінюються }
\rho_m,\quad
\vect v,\quad
p,\quad
\vect B?
\end{equation*}

У~кінетичній теорії запитання стає детальнішим:

\begin{equation*}
\text{як змінюється }
f_s(\vect r,\vect v,t)
\text{ для частинок різних швидкостей?}
\end{equation*}

Перший опис містить менше інформації, але тому є простішим. Другий містить значно
більше інформації і~дозволяє досліджувати ефекти, які МГД принципово не може побачити.

% --------------------------------------------------------
\subsubsection{Коли МГД перестає бути достатньою?}
% --------------------------------------------------------

Для того щоб МГД-опис був виправданим, хвиля повинна бути достатньо повільною та
великомасштабною порівняно з~характерними мікроскопічними процесами плазми. Один із
важливих критеріїв для магнітованої плазми має вигляд

\begin{equation}\label{eq:mhd-wave-kinetic-scale}
k_\perp r_L\ll1,
\end{equation}

де $r_L$ --- характерний ларморівський радіус частинок.

У~цій межі хвиля не розрізняє детальної структури ларморівських орбіт: для неї рух
частинок може бути описаний через усереднені макроскопічні величини.

Якщо ж

\begin{equation*}
k_\perp r_L\sim1,
\end{equation*}

структура хвилі на поперечному масштабі стає порівнянною з~ларморівським рухом частинок.
Тоді усереднення, яке лежить в~основі простої МГД, вже не приховує всієї важливої фізики.

Аналогічна ситуація виникає, коли частота хвилі наближається до характерних частот
мікроскопічного руху. Наприклад, для магнітованої плазми істотними можуть стати
циклотрононні резонанси, якщо

\begin{equation*}
\omega\sim\omega_c.
\end{equation*}

Нарешті, навіть у~великомасштабній плазмі МГД може виявитися недостатньою, якщо важливою
стає залежність відповіді від швидкості частинок. Саме тоді з'являються кінетичні
резонанси, які не можна описати лише густиною, середньою швидкістю та тиском.

Отже, межа між МГД та кінетикою визначається не просто словом `велика'' або `мала''
довжина хвилі. Вона визначається порівнянням характерних масштабів і~частот задачі
з~мікроскопічними масштабами руху частинок.

% --------------------------------------------------------
\subsubsection{Що змінюється при переході до кінетики?}
% --------------------------------------------------------

Це дозволяє сформулювати головну ідею наступного розділу.

У~МГД ми вже знаємо, що плазма може підтримувати хвилі. Ми знаємо їхні швидкості,
поляризації та взаємодію між магнітним полем і~рухом середовища.

Тепер нас цікавить, \emph{що відбувається всередині цієї макроскопічної картини}.

Які частинки забезпечують хвильовий рух? Чи всі частинки реагують на хвилю однаково?
Що відбувається з~частинками, швидкість яких близька до фазової швидкості хвилі?
Чи може хвиля передавати їм свою енергію? Чи може, навпаки, отримувати від них енергію?

На ці питання МГД відповісти не може, оскільки в~ній немає самої змінної $\vect v$ як
незалежного аргументу функції стану.

Кінетичний опис повертає цю інформацію. Ми знову розглядаємо

\begin{equation}\label{eq:mhd-to-kinetic-distribution}
f_s=f_{0s}+\delta f_s,
\end{equation}

але тепер $\delta f_s$ залежить від $\vect v$. Це дозволяє простежити не лише за
колективним рухом плазми, а~й за тим, як окремі групи частинок обмінюються енергією
з~хвилею.

Таким чином, перехід від МГД до кінетики є переходом

\begin{equation*}
\tcbhighmath{
\text{від руху середовища}
\quad\longrightarrow\quad
\text{до руху частинок, що утворюють це середовище}.
}
\end{equation*}

МГД-хвилі — це граничний випадок ширшого кінетичного опису. Щоб зрозуміти, звідки береться ця межа і що відбувається за її межами, тепер потрібно перейти від макроскопічних величин до функції розподілу окремих частинок. Саме цей перехід ми тепер і~здійснимо в наступному розділі.


%% ========================================================
\section*{Підсумки}
%% ========================================================



\begin{itemize}
\item Хвиля виникає тоді, коли збурення породжує відновлювальну силу та може передаватися
  від однієї ділянки середовища до іншої.

\item У~плазмі відновлювальну дію можуть створювати тиск плазми, магнітний натяг та
  магнітний тиск.

\item Звукова хвиля визначається переважно тиском і~поширюється зі швидкістю $c_s$.

\item Альвенівська хвиля є поперечним збуренням магнітних силових ліній; її швидкість
  визначається $v_A$.

\item У~магнітозвукових хвилях тиск плазми та магнітне поле діють одночасно, тому
  виникають швидка та повільна моди.

\item МГД описує лише частину спектра плазмових хвиль. На малих масштабах та поблизу
  характерних частот руху частинок потрібен кінетичний опис.
\end{itemize}



%% --------------------------------------------------------
\section*{Контрольні питання}
%% --------------------------------------------------------



\begin{Questions}
\begin{quizlist}

\item Що таке відновлювальна сила? Чому без неї не виникає хвильове поширення збурення?

\item Чим фізично відрізняються звукова та альвенівська хвилі?

\item Чому швидкість Альвена залежить від $B$ та $\rho_m$, але не визначається
  температурою плазми?

\item Яка роль магнітного тиску у~магнітозвуковій хвилі?

\item Чому при довільному напрямку поширення виникають дві магнітозвукові моди?

\item Що відбувається з~МГД-описом, коли довжина хвилі стає порівнянною з~мікроскопічними
  масштабами плазми?

\item Чому кінетична теорія потрібна не для \enquote{кращого} опису взагалі, а~лише для
  процесів, у~яких стає важливою структура функції розподілу?

\end{quizlist}
\end{Questions}

